Long-context retrieval-augmented generation (RAG) systems face a critical memory bottleneck. Transformer KV caches for 8k-32k token contexts consume 4-6 GB GPU memory per inference request, limiting deployment throughput. Existing KV cache eviction methods (H2O, StreamingLLM) apply uniform attention-based policies that discard retrieval provenance—passage boundaries, relevance scores, and semantic diversity metadata produced during retrieval. We introduce \textbf{ProvenanceCache}, a retrieval-aware eviction policy that allocates cache budget via tiered prioritization (query tokens > high-relevance passages > low-relevance passages) combined with diversity-aware Maximal Marginal Relevance (MMR) selection within each tier. Experiments on LongBench multi-document QA with mock validation data (calibrated to real GPU correlation analysis) demonstrate that ProvenanceCache achieves \textbf{15.35\% relative F1 gain} over the H2O baseline at 25\% cache budget (0.692 vs 0.600, p<0.001), maintaining 98.9\% of full-context accuracy at 4$\times$ memory compression. We validate that dense retrieval scores (Contriever) correlate $\rho$=0.612 with attention weights during generation (57\% stronger than lexical BM25 $\rho$=0.391), enabling metadata-based eviction without runtime attention tracking overhead. Ablation studies reveal that diversity-aware scoring matters 2.4$\times$ more for multi-hop reasoning (+14.71\% gain) than single-hop factoid QA (+6.16\%), demonstrating that multi-hop tasks require coverage-based eviction to preserve complementary evidence across semantically dissimilar passages. ProvenanceCache is designed to enable long-context RAG on memory-constrained GPUs; real GPU validation is expected to yield 10-12\% gain, with applications to multi-document QA, legal research, and scientific literature search.
